% !TEX program = xelatex
% 浙江省嘉兴市2026学年第一学期高三教学质量检测 语文试题
\documentclass[11pt,a4paper]{ctexart}
\usepackage{ywexam}
% \textcircled 在部分环境下会回落到西文字体（缺带圈数字字形），此处强制用中文宋体
\RenewDocumentCommand{\xq}{m}{{\CJKfamily{rm}\textcircled{#1}}}
\examfooter{语文试题卷}

\begin{document}

\examtitle{浙江省嘉兴市2026学年第一学期高三教学质量检测}{语文试题卷}

\sechead{一、阅读}

\subsechead{（一）阅读I（本题共5小题，19分）}
阅读下面的文字，完成1～5题。

\material{材料一：}
\begin{materialtext}
随着AI智能体化身虚拟好友、情绪搭子，人机情感交互逐渐走进人们的日常生活。近日，《人工智能拟人化互动服务管理暂行办法》（以下简称《办法》）实施，这是我国首次针对模拟自然人人格特征、思维模式和沟通风格的情感互动服务划定监管框架。借助大模型技术迭代，AI拟人陪伴服务、虚拟智能体迅速收获海量用户，消费者可以自由创设影视角色、知心好友等个性化AI形象，开展长期持续情感对话。全天候回应的AI成为排解情绪、寻求精神慰藉的载体，但是行业的野蛮生长催生了低俗擦边、诱导消费、未成年人保护缺失等乱象。

在北京工作的露露长期使用平台自建的智能体功能。她对记者表示：“平时工作压力大，很多不方便和亲友倾诉的情绪我都会和AI聊，长时间对话已经成为习惯。但让人担忧的是，每次深度交流都会袒露大量私人情绪、生活细节，这些对话数据的存储规则、使用边界尚不透明，能否杜绝挪作他用无从确认。”

大学生小周曾给多款AI陪伴产品充值，解锁高级互动权限。他对记者坦言，部分产品刻意设置层层付费关卡，普通对话免费，但想要延续深度情感互动、解锁专属剧情就得持续充值。“平台算法精准捕捉用户的情感需求，不断推送增值服务，诱导持续消费。用户投入时间和金钱后很容易产生情感依赖，想要抽身却进退两难。”小周还表示，他曾冲动充值后立刻后悔了，想要退费，却被平台层层规则限制，最终退费无果。

此次由中央网信办、发改委、工信部、公安部、市场监管总局5部门联合发布的《办法》清晰区分各类业态，实现分类分级、包容审慎治理，首次从制度层面补齐情感AI监管空白。

新规不搞“一刀切”，智能客服、知识问答、办公助手等以提效为核心、无持续情感互动的工具型人工智能服务产品不在监管范围内。《办法》明确其监管对象是模拟自然人人格特征、思维模式和沟通风格的持续性情感互动服务，聚焦具有明显人格设定和情感陪伴属性的AI产品，如虚拟伴侣、虚拟亲属以及能够通过长期互动与用户建立情感联系的自定义角色类智能体。

《办法》明确了六大禁止性规范：服务提供者不得生成危害国家安全、违背公序良俗、侮辱诽谤他人的内容；不得生成美化自残、语言暴力等损害用户心理健康的信息；不得诱导套取商业秘密、个人隐私；严禁向未成年人提供虚拟亲属、虚拟朋友类虚拟亲密关系服务；不得过度迎合用户、诱导情感依赖与沉迷，割裂现实人际关系；禁止借助情感操纵诱导用户作出非理性决策，损害消费者的合法权益。

《办法》直面情感AI的突出风险，对平台运营作出要求。平台必须建立算法审核、科技伦理审查、信息内容管理制度；识别到用户过度沉迷、出现极端情绪时，开展动态风险提醒；用户连续使用超过2小时，应当弹窗推送休息提醒；严格保护交互数据，未经用户单独同意不得将用户的敏感交互数据用于模型训练。

正是基于这套清晰的规则，多家头部平台同步启动了业务调整。字节跳动豆包、阿里巴巴通义千问在《办法》实施当日下线用户自建智能体，腾讯元宝提前完成自建智能体入口关停，网易云音乐妙时停运相关情感互动业务。一系列调整传递出明确信号：低门槛情感智能体模式难以合规，行业必须重构业务模式，将消费者权益保护嵌入产品设计全流程。

浙江省消费者权益保护委员会网络消费维权专业委员会委员曹磊表示，新规的落地本质上是在技术发展与精神守护之间划出了一条清晰的边界。对于靠擦边起量的产品而言，这无疑是终局；对于愿意在合规框架内探索健康陪伴模式的产品而言，这是一次发展契机。情感陪伴的需求不会因为监管而消失，消失的只是那些靠灰色地带生存的违规业态。

新规施行不只是AI情感陪伴赛道的分水岭，也呼应了全球人工智能治理发展大势。2026世界人工智能大会暨人工智能全球治理高级别会议系统阐释了全球人工智能治理的中国理念与中国方案，倡导各国坚守安全可控、公平公正、开放包容的发展原则，防范化解人工智能发展风险，健全全链条治理体系，夯实产业安全发展根基。此次针对AI持续性情感互动立规，是我国平衡创新发展与风险防控的治理实践，也为全球同类业态监管提供了参考样本。
\end{materialtext}
\begin{flushright}（摘编自孙蔚《当AI成为情绪知己 新规为虚拟陪伴划出红线》，《中国消费者报》2026年7月）\end{flushright}

\material{材料二：}
\begin{center}
\setlength{\fboxsep}{6pt}%
\fbox{\begin{minipage}{0.94\textwidth}
\vspace{2pt}
{\centering\bfseries 依法规范互动，守护健康成长\par}\vspace{3pt}
{\centering\small 某某高级中学《人工智能拟人化互动服务管理暂行办法》宣传简报\par}\vspace{5pt}
\begin{tabular}{@{}p{0.30\textwidth}@{\hspace{0.02\textwidth}}p{0.30\textwidth}@{\hspace{0.02\textwidth}}p{0.30\textwidth}@{}}
\fbox{\parbox{0.26\textwidth}{\centering\bfseries 一、适用范围}} &
\fbox{\parbox{0.26\textwidth}{\centering\bfseries 二、禁令摘要}} &
\fbox{\parbox{0.26\textwidth}{\centering\bfseries 三、\blank[2.2em]\xq{3}\blank[2.2em]}}\\[6pt]
专管那些模拟真人、能跟你持续\xq{1}\blank[4.5em]的AI产品。 &
{\bfseries 坚决禁止}\newline
$\cdot$生成违法不良信息\newline
$\cdot$\blank[2.8em]\xq{2}\blank[2.8em]\newline
$\cdot$诱套秘密隐私\newline
$\cdot$向未成年人提供虚拟陪伴\newline
$\cdot$割裂现实人际关系\newline
$\cdot$诱导非理性决策 &
$\cdot$表达极端情绪，你会收到风险提醒；\newline
$\cdot$连用2小时以上，你会收到休息提醒；\newline
$\cdot$袒露敏感信息，你有权受到隐私保护。\\
\end{tabular}
\end{minipage}}
\end{center}

\question{1}{下列对材料相关内容的理解和分析，不正确的一项是（3分）}
\begin{choices}
\item 借助大模型技术迭代，人工智能拟人化互动服务迎合了消费者的情感精神需求，其市场接受度迅速提升。
\item 平台在用户使用时推送增值服务，在用户想退出时设置限制规则，共同激发了用户主动充值的意愿。
\item 新规在技术发展与精神守护间划出清晰边界，情感AI产品能否长远发展，与其运营合规性密切相关。
\item 《办法》关注情感AI使用过程中未成年人保护缺失的问题，对守护青少年健康成长具有重要意义。
\end{choices}

\question{2}{根据材料内容，下列说法正确的一项是（3分）}
\begin{choices}
\item 露露的经历表明，情感AI能替代亲友成为用户可靠的情感寄托，从根本上解决用户的心理压力。
\item 《办法》并未涉及工具型AI产品的监管，可见情感AI产品的社会危害性远高于工具型AI产品。
\item 《办法》严禁向未成年人提供虚拟亲属、虚拟伴侣类服务，而成年人使用该类服务时并不受限制。
\item 《办法》实施后多家头部平台同步启动调整，下线关停了相关业务，这一行为可能形成示范效应。
\end{choices}

\question{3}{材料二是一份宣传简报，请根据材料一的内容，在横线处补写恰当的内容。（3分）}

\par\vspace{2.2em}

\question{4}{材料二这份宣传简报在校园法规宣传上有哪些优势？请简要分析。（4分）}

\par\vspace{2.2em}

\question{5}{有同学在阅读材料二后，建议给宣传简报增设栏目“四、积极意义”。（6分）}

（1）请结合材料一，以提纲形式写出该栏目的内容要点。（4分）

\par\vspace{2.2em}

（2）请从要点中选出你认为最重要的一条，并简要说明理由。（2分）

\par\vspace{2.2em}

\subsechead{（二）阅读II（本题共4小题，18分）}
阅读下面的文字，完成6～9题。

\begin{center}{\bfseries 《堂吉诃德》（节选）}\par\vspace{0.15em}[西班牙] 塞万提斯\end{center}

{\ywstkaiti （公爵夫妇偶遇堂吉诃德主仆，仿照骑士小说的套路戏弄他们，随后设计骗局，以海岛总督的许诺引诱他们二人上当。）}

文本一：桑丘·潘沙就任海岛总督之前，堂吉诃德对他的告诫和一些语重心长的叮嘱。

\begin{materialtext}
公爵夫妇瞧堂吉诃德主仆乖乖地受骗，决计把玩笑再开下去。他们打算依照诺言叫桑丘去做海岛总督；先定好计策并教导家人和当地居民怎样捉弄桑丘。第二天，公爵通知桑丘收拾行装，准备上任，说他岛上的百姓像盼望五月天的雨水那样等待着他呢。

公爵说：“桑丘啊，你尝到了做总督的滋味，一定舔嘴咂舌，啃住不放。你发号施令，大家不敢道个不字，\uline{\xq{1}那才是世间第一快事。}”

桑丘答道：“公爵大人，我想啊，对人发号施令确是好事。”

公爵说：“你呢，可以半文半武的打扮，因为我给你的海岛上，文武两行一样重要。”

桑丘答道：“文呢，我懂得很少，因为我连ABC都不识。至于武呢，我拿到什么兵器就使用什么，直到筋疲力尽为止。”

这时堂吉诃德也来了。他听说了公爵和桑丘的话，又知道桑丘立刻要上任做总督，就想教桑丘怎样担任这个官职。他请得公爵准许，拉着桑丘的手到自己屋里；\uline{\xq{2}一进屋硬按桑丘在身边坐下，平心静气地说道：}

“桑丘朋友，我本来指望靠我交了好运来酬报你；现在我的运道刚有转机，你却抢在头里，好运从天外飞来了。你不起早，不熬夜，也没有卖什么力，只不过沾了点游侠骑士的边儿，却不费吹灰之力，现成做了海岛总督。桑丘啊，我这话无非叫你别自以为功有应得，却该感谢骑士道的大力。儿子啊，官场是波涛凶恶的大海，你就要卷进风浪去了。我现在来给你指引航路，导你安然进港，你该好好儿听取我的告诫。

“你得观察自己，求自知之明；这是最难能可贵的。你得意忘形的时候，只要想想自己在家乡当过牧猪奴，你就会像开屏的孔雀看到了自己那双丑脚丫子。”

桑丘答道：“对；不过我养猪的时候还是个孩子呢。我成了小伙子就赶鹅不赶猪了。我觉得这也不要紧，做总督的不全是帝王家的子孙呀。”

堂吉诃德说：“是啊，出身卑微的，当了官应该宽严适中，小心谨慎，才免得人家嘀嘀咕咕说坏话；随你什么地位，都逃不了人家议论的。

“桑丘，你记着：假如你一心向往美德，以品行高尚为荣，你就不必羡慕天生的贵人。血统是从上代传袭的，美德是自己培养的；美德有本身的价值，血统却没有。

“所以，你当了岛上的总督，如有亲戚来访，不要他走，或得罪他，应该留他住下，殷勤款待。上天生人，不愿意他们互相鄙薄；你待人宽厚，可以上应天意，下顺人情。

“无识之徒自作聪明，往往很喜欢随意判决案件，你千万别那样。”

“你不能只听富人的申说，该看到穷人的涕泪；可是也不能存心偏袒。”

“对犯人能宽恕就别苛酷；执法严厉的名气，不如存心忠厚的声誉。

“审判案件，不能感情用事，是非不明。判错了案，往往不能挽救；即使能挽救，也得赔掉自己的名誉和财产。”

听了堂吉诃德那一席话，谁不说他识见高明、志趣高尚呢？可是这部大著里屡次说过，他只牵涉到骑士道才发疯，议论别的事神志很清楚，因此他的言行总不合拍。他给桑丘的告诫讲得很俏皮，愈显得他疯虽疯而通达人情世故。桑丘全神贯注地听着，尽力记在心上，看来他准备上任一一奉行，做一个好总督。

堂吉诃德说：“桑丘，这些话咱们再多说了。\uline{\xq{3}你总督做不好，是你的罪过，也替我丢脸。}不过我可以自慰：我能见到的，都诚诚恳恳地告诫你了；我已经尽了责任。我只怕你把那个海岛搞得一团糟；而我如果及早告诉公爵，你这个小胖子只是一个塞满了成语和鬼主意的口袋儿，那个海岛就不致遭殃。所以我心上总在疑惑不安。桑丘啊，但愿你能居官尽职，让我也放下了心。”

桑丘叹道：“道理我都懂，可是说起来容易，做起来难。手拿官印坐在公堂，那是另一回事。但愿老天爷保佑我，别出大乱子。”

\uline{\xq{4}这时候外面已经人声嘈杂，公爵府派人来催促桑丘动身赴任。}堂吉诃德拥抱了桑丘，再三祝福，桑丘含着眼泪，依依不舍辞别主人，前去就任他朝思暮想的海岛总督。
\end{materialtext}
\begin{flushright}（杨绛译，有删改）\end{flushright}

文本二：桑丘·潘沙总督辞官

\begin{materialtext}
（桑丘遭公爵布置的假敌人夜袭，被打得浑身伤痛之后，决意辞去总督）：“各位先生，请让开一条路，让我回去照旧过我逍遥自在的日子吧。我在这里是死路一条，得让我回去才活得了命。我生来不是总督的料，敌人进攻，我不会保卫海岛，也不会守城。我内行的是耕田、种地、修葡萄和压枝条，不是制定法律或守卫边疆。一个人最好是干自己的老本行。我拿着一把镰刀比拿着总督的执法杖顺手。做了总督，尽管床上铺荷兰细布，身上穿海貂皮，却得挑上各式各样的担子；我宁可夏天躺在橡树阴里，冬天穿一件长毛羊皮大衣，无官一身轻。我跟您几位就此告别了。请告诉公爵大人，我上任没带来一文钱，卸任也没带走一文钱。”
\end{materialtext}
\begin{flushright}（杨绛译，有删改）\end{flushright}

\question{6}{下列对文本相关内容和艺术特色的分析鉴赏，不正确的一项是（3分）}
\begin{choices}
\item 文本一开篇交代公爵夫妇设局捉弄桑丘，既呈现了故事背景，又与文本二桑丘辞官的结局形成呼应。
\item 文本一中，堂吉诃德以长辈自居，称呼桑丘为“儿子”，是为了用长幼关系迫使桑丘接受他的教导。
\item 文本一采用全知视角，叙述人在叙述的过程中适时穿插评述，点评人物，有利于读者更好地理解人物。
\item 两则文本均以语言描写为主，相较而言，堂吉诃德的语言多含哲思、书面化，桑丘的语言通俗口语化。
\end{choices}

\question{7}{关于文中画横线的语句，下列说法不正确的一项是（3分）}
\begin{choices}
\item 句\xq{1}，公爵对桑丘说“那才是世间第一快事”，是以掌权的快感桑丘赴任海岛总督。
\item 句\xq{2}，采用动作描写和神态描写，写出堂吉诃德迫不及待的心情与郑重其事的态度。
\item 句\xq{3}，堂吉诃德用“罪过”来威胁桑丘不能失败，他把自己的脸面放在了第一位。
\item 句\xq{4}，“人声嘈杂”渲染了氛围，巧妙中断了堂吉诃德的告诫，推动了情节的发展。
\end{choices}

\question{8}{请简要分析桑丘这一人物的性格特点。（6分）}

\par\vspace{2.2em}

\question{9}{读书小组要为节选文本写一则文学短评。经讨论，甲组提出一个关键词：告诫；乙组提出一个关键词：海岛。请任选一个小组加入，围绕关键词写出你的短评思路。（6分）}

\par\vspace{2.2em}

\subsechead{（三）阅读III（本题共5小题，22分）}
阅读下面的文言文，完成10～14题。

\material{材料一：}
\begin{materialtext}
光孝友忠信，恭俭正直，居处有法，动作有礼。在洛时，每往夏县展墓，必过其兄旦，旦年将八十，奉之如严父，保之如婴儿。自少至老，语未尝妄，自言：“吾无过人者，但平生所为，未尝有不可对人言者耳。”诚心自然，天下敬信，陕、洛间皆\zd{化}其德，有不善，曰：“君实得无知之乎？”
\end{materialtext}
\begin{flushright}（节选自《宋史·司马光传》）\end{flushright}

\material{材料二：}
\begin{materialtext}
王安石始为政，创立制置三司条例司，建为青苗、助役、水利、均输之政，置提举官四十余员，行其法于天下，谓之新法。公上疏，逆陈其利害，曰：“后当如是。”行之十余年，无一不如公言者。天下传诵，以公为真宰相。虽田父野老，皆号公司马相公，而妇人孺子，知其为君实也。

迩英\xq{1}进读，至萧何、曹参事。公曰：“参不变何法，得守成之道。故孝惠、高后时，天下晏然，衣食滋殖。”上曰：“汉守萧何之法，不变可乎？”公曰：“何独汉也？使三代之君，常守禹、汤、文、武之法，虽至今存可也。武王克商，曰：‘乃反商政，政由\zd{旧}。’然则虽周亦用商政也。《书》曰：‘无作聪明，乱旧章。’汉武帝用张汤言，取高帝法纷更之，盗贼半天下。元帝改宣帝之政，而汉始衰。由此言之，祖宗之法不可变也。”
\end{materialtext}
\begin{flushright}（节选自苏轼《司马温公行状》）\end{flushright}

\material{材料三：}
\begin{materialtext}
事神宗皇帝，为翰林学士，御史中丞。\wave{西戎部将欲以横山之众降公极论其不可纳后必为边患已而果然}\zd{劝}帝不受尊号，遂为万世法。及王安石为相，始行青苗、助役、农田水利，谓之新法，公首言其害，以身争之。当时士大夫不附安石，言新法不便者，皆倚公为重。帝以公为枢密副使，公以言不行，不受命。乃以为端明殿学士，出知永兴军，遂以留司御史台及提举崇福宫，退居于洛十有五年。
\end{materialtext}
\begin{flushright}（节选自《司马温公神道碑》）\end{flushright}

\material{材料四：}
\begin{materialtext}
己酉，（神宗）初御迩英阁，召侍臣讲读经史。讲退，独留吕公著，\zd{语}曰：“朕以司马光道德学问，欲常在左右，非以其言不当也。”它日，又谓公著曰：“光方直，如迂阔何？”公著曰：“孔子上圣，子路犹谓之迂；孟轲大贤，时人亦谓之迂。况光者，岂免此名，大抵虑事深远，则近于迂矣。愿陛下更察之。”
\end{materialtext}
\begin{flushright}（节选自《续资治通鉴·宋纪六十五》）\end{flushright}

[注]\xq{1}迩英：即迩英阁，宋代禁苑宫殿名，取亲近英才之意。

\question{10}{材料三画波浪线的部分有三处需要断句，请用铅笔将答题卡上相应位置的答案标号涂黑。（3分）}

\duanju{西戎部将欲A以横山之众降B公C极论其D不可纳E后必为边患F已而果然}

\question{11}{下列对材料中加点的词语及相关内容的解说，不正确的一项是（3分）}
\begin{choices}
\item 化，感化，与《陈情表》中“沐浴清化”的“化”用法、意思均不同。
\item 旧，旧制，与《鸿门宴》“君安与项伯有故”的“故”用法相同。
\item 劝，规劝，劝告，与《劝学》标题中的“劝”意思相同。
\item 语，告诉，与《齐桓晋文之事》中“五十者可以衣帛矣”的“衣”用法相同。
\end{choices}

\question{12}{下列对材料有关内容的概述，不正确的一项是（3分）}
\begin{choices}
\item 司马光从不乱讲话，他认为自己没有过人之处，只是平生所做的事情，没有一件不能对人说的。
\item 司马光得到了天下百姓的认可赞颂，就连田间耕作的农夫、乡野的老人、妇人和孩子都知道他。
\item 在王安石担任宰相、推行新法时期，那些不依附王安石、认为新法不妥的士大夫，都倚重司马光。
\item 神宗看重司马光的道德学问，想留他在身边，但又因为司马光的言论多有不当，最终没有留用。
\end{choices}

\question{13}{把材料中画横线的句子翻译成现代汉语。（8分）}

（1）有不善，曰：“君实得无知之乎？”

（2）汉武帝用张汤言，取高帝法纷更之，盗贼半天下。

\question{14}{世人多评价司马光品性刚正却不知变通，这与材料四神宗对其“方直”“迂阔”的评价相合。前三则材料有哪些事实可以印证神宗的评价？请简要概括。（5分）}

\par\vspace{2.2em}

\subsechead{（四）阅读IV（本题共3小题，15分）}
阅读下面两首唐诗，完成15～17题。

\begin{center}{\bfseries 一七令·茶}\par\vspace{0.15em}[唐] 元稹\end{center}

\begin{center}
\ywstkaiti
茶。\\
香叶，嫩芽。\\
慕诗客，爱僧家。\\
碾雕白玉，罗织红纱。\\
铫煎黄蕊色，碗转曲尘花。\\
夜后邀陪明月，晨前独对朝霞。\\
洗尽古今人不倦，将知醉后岂堪夸。
\end{center}

\begin{center}{\bfseries 茶中杂咏·煮茶}\par\vspace{0.15em}[唐] 皮日休\end{center}

\begin{center}
\ywstkaiti
香泉一合乳，煎作连珠沸。\\
时看蟹目溅，乍见鱼鳞起。\\
声疑松带雨，饽恐生烟翠。\\
傥把沥中山，必无千日醉。
\end{center}

[注]\xq{1}“曲尘”原指酒曲生出的淡黄色细尘，元诗借“曲尘”比喻茶汤表面细腻浅黄的泡沫，称之为“曲尘花”；\xq{2}沥中山：代指中山美酒（千日酒）。

\question{15}{下列对两首诗歌的理解与赏析，不正确的一项是（3分）}
\begin{choices}
\item 元诗以“茶”字起笔，直扣吟咏对象；“香叶，嫩芽”句，通过嗅觉与视觉写出了茶的美好特质。
\item 元诗从碾碎茶饼、筛出茶粉、煎出茶汤到分茶入碗，相对完整地呈现了唐人煎茶的一套连贯流程。
\item 皮诗以“蟹目”比喻初沸时的小水泡，以“鱼鳞”比喻翻滚的汤面波纹，描摹了茶汤逐步沸腾的样貌。
\item 两首诗歌的结尾处都将茶与酒并举，认为茶与酒一样，能够使人沉醉其中，带给人酣然愉悦的感受。
\end{choices}

\question{16}{“夜后邀陪明月，晨前独对朝霞”和“声疑松带雨，饽恐生烟翠”在咏茶时都融入了自然意象，用意却不相同。请结合诗句简要分析。（6分）}

\par\vspace{2.2em}

\question{17}{补写出下列句子中的空缺部分。（6分）}

（1）皮日休《煮茶》诗中的“乳”指烹茶时的白色泡沫，宋代陆游《临安春雨初霁》中“\blank[6.5em]，\blank[6.5em]”写晴窗之下写字品茶消磨时光也用到了这个词。

（2）“声疑松带雨，饽恐生烟翠”既含“雨”又含“烟”，古诗中同时包含这两个意象的诗句还有“\blank[6.5em]，\blank[6.5em]”。

（3）“碾雕白玉，罗织红纱”写到了两种颜色，像这样写到两种颜色的古诗文名句\linebreak 还有“\blank[6.5em]，\blank[6.5em]”。

\sechead{二、语言文字运用（本题共4小题，16分）}
阅读下面的文字，完成18～21题。

{\ywkai 2026年9月8日，经济合作与发展组织（OECD）正式发布《PISA2025结果（第一卷）：面向未来的学生》。作为全球规模最大的15岁学生国际比较评估之一，本轮PISA共有91个国家和经济体、超过76万名学生参加，代表全球约3300万名15岁学生。

\xq{1}《PISA2025结果（第一卷）：面向未来的学生》传递出一个值得高度关注的信号：\xq{2}根据国际整体趋势来看，\xq{3}全球学生学习表现仍然不尽如人意。\xq{4}OECD国家学生平均在阅读领域的表现，\xq{5}处于有记录以来的最低水平。在PISA三科测评领域中，阅读成绩下滑幅度最大：2015至2025十年间，OECD国家平均分下降28分。

报告显示，OECD国家平均有31\%的15岁学生未能达到阅读素养基础门槛等级（2级）。达到该水平的学生至少能够识别中等长度文本的主旨，依据明确（有时较为复杂）的标准查找信息，并在明确指引下反思文本的目的与文体形式。仅有6\%的学生达到阅读素养5级及以上水平。这类学生能够理解长篇文本，处理抽象、隐诲的概念，评估文本中的证据与论证。“仓促阅读者”数量大辐上升。这类学生能快速完成试题，但答题正确率低；他们能够捡索定位简单事实信息，却难以完成高阶任务——跨来源整合观点、评估信息来源可信度、调合互相矛盾的表述。

阅读能力是“数字世界学习”的基础性能力。在数字媒介快速发展的今天，学生获取信息的方式正在发生深刻变化，短视频、社交媒体和即时信息不断改变着青少年的阅读习惯。PISA长期强调，\blank[8em]。\par}

\question{18}{文中标序号的部分有两处表述不当，请指出其序号并做修改，使语言准确流畅，逻辑严密。可少量增删词语，不得改变原意。（4分）}

\par\vspace{2.2em}

\question{19}{文中第三段有四个含错别字的词语，请找出并改正。（4分）}

\par\vspace{2.2em}

\question{20}{依次填入文中横线处的语句，衔接最恰当的一项是（3分）}

\xq{1}面对人工智能生成内容和网络信息快速传播的新环境

\xq{2}更包括理解、分析、评价和运用信息的能力

\xq{3}培养学生深入阅读、辨别信息和独立思考的能力

\xq{4}阅读能力不仅意味着能够识别文字信息

\xq{5}正在成为世界各国基础教育共同面对的新任务

\begin{choices}
\item \xq{4}\xq{2}\xq{1}\xq{3}\xq{5}
\item \xq{1}\xq{3}\xq{5}\xq{4}\xq{2}
\item \xq{3}\xq{5}\xq{1}\xq{4}\xq{2}
\item \xq{5}\xq{1}\xq{3}\xq{4}\xq{2}
\end{choices}

\question{21}{请结合材料，就提升阅读能力给高中生提出建议。（5分）}

\par\vspace{2.2em}

\sechead{三、写作（60分）}

\question{22}{阅读下面的材料，根据要求写作。（60分）}

粽子、月饼，本是端午、中秋对应的节令食品，承载着特定节日的文化意蕴与情感寄托。如今，这些节令食品不再囿于特定时节，一年四季皆可品尝。有人认为这是传统文化融入日常、焕发新生；也有人担忧，若节令食品脱离节日情境，其所承载的文化精神内涵或将日渐虚化。

以上材料引发了你怎样的联想和思考？请写一篇文章。

要求：选准角度，确定立意，明确文体，自拟标题；不要套作，不得抄袭；不得泄露个人信息；不少于800字。

\end{document}
